%% kapa: exam=1 round=34 year=2023 period=2 form=A subject=appraisal_law subject_ko=감정평가관계법규 seq=36 number=36
%% source: Q-Net 감정평가사 기출문제 (cst00309) artlSeq=5211789 file=raw/1008/5211789/2023년도 제34회 감정평가사 1차시험 2교시.pdf
%% conversion: kordoc 4.24.51 -> tools/gen1.py (auto); review=proofread-sonnet; flags=none
\begin{question}{36}
도시 및 주거환경정비법령상 조합에 관한 설명으로 옳은 것은?
\begin{choices}
  \item 조합이 정관의 기재사항인 조합임원의 수를 변경하려는 때에는 시장·군수등의 인가를 받아야 한다.
  \item 건설되는 건축물의 설계 개요를 변경하려면 조합총회의 의결을 거쳐야 한다.
  \item 토지등소유자의 수가 100인을 초과하는 경우 조합에 두는 이사의 수는 7명 이상으로 한다.
  \item 조합임원의 임기는 5년 이하의 범위에서 정관으로 정하되, 연임할 수 없다.
  \item 조합의 대의원회는 조합원의 10분의 1 이상으로 구성하며, 조합장이 아닌 조합임원도 대의원이 될 수 있다.
\end{choices}
\end{question}
\answerall
\answernote{정답표에 1,2,3,4,5 모두 정답으로 표기(전원 정답)}
